\section*{前言}
\phantomsection
\label{frontmatter:preface}
\addcontentsline{toc}{section}{前言}

本专著是先前一本法文著作\MBUnresolvedReference{bib:source-24}{[24]}的修订扩充版, 该书出版于 2006 年. 我们有两个目标. 第一, 向读者介绍流体力学中的建模与数学分析. 我们处理的核心模型是不可压缩 Stokes 方程和 Navier--Stokes 方程, 它们的推导将在\MBUnresolvedReference{chap:fluid-mechanics-equations}{Chapter~I}给出. 第二, 以足够一般的方式介绍数学工具, 使读者能够用这些工具研究许多其他类型的非线性演化偏微分方程.

本着这一精神, 我们力求使本书尽可能自成体系. 开始阅读所需的先修知识仅为微积分, 积分和泛函分析中的基本结果. \MBUnresolvedReference{chap:analysis-tools}{Chapter~II}回顾本书所用的一些较为标准的分析结果, 其中包括 Banach 空间中弱收敛与弱-* 收敛的定义和性质, 分布理论简介, 基本紧性结果, Bochner 积分及相关空间, 谱理论的基本结果等.

在\MBUnresolvedReference{chap:sobolev-spaces}{Chapter~III}中, 我们介绍并研究 $\mathbb{R}^d$ 中 Lipschitz 区域上的 Sobolev 空间. \MBUnresolvedReference{sec:ch03-domains}{Section~1}包含这类区域的主要定义和性质. 特别地, 我们说明如何定义这类区域边界上所定义函数的积分, 并证明基本的 Stokes 公式. 随后在\MBUnresolvedReference{sec:ch03-sobolev-spaces-on-lipschitz-domains}{Section~2}中介绍并深入研究 Sobolev 空间. 适当的磨光算子在其中发挥重要作用, 不仅在本章如此, 在后续各章亦然; 借助这些算子, 我们可以证明各种函数空间中光滑函数集合的稠密性. 我们还研究嵌入定理, 延拓算子, 迹算子与迹提升算子, Sobolev 空间的对偶理论, 以及十分重要的 Poincaré 不等式和 Hardy 不等式. \MBUnresolvedReference{sec:ch03-calculus-near-boundary}{Section~3}介绍充分光滑区域边界附近适当的法向/切向坐标. 这一框架可用于研究切向 Sobolev 空间, 从而证明椭圆问题解直至边界的正则性, 例如 Stokes 问题. \MBUnresolvedReference{sec:ch03-laplace-problem}{Section~4}通过完整研究带 Dirichlet 或 Neumann 边界条件的 Laplace 问题来说明上述所有概念.

\MBUnresolvedReference{chap:steady-stokes-equations}{Chapter~IV}讨论稳态(不可压缩)Stokes 方程. 对任何不可压缩流体力学模型的研究而言, 这都是一个核心主题. 该理论的一个基本工具是 Nečas 不等式, \MBUnresolvedReference{sec:ch04-necas-inequality}{Section~1}给出它的完整证明. \MBUnresolvedReference{sec:ch04-gradient-fields}{Section~2}利用此不等式研究梯度场与无散向量场之间的关系, 这对于构造 Stokes 系统中的压力至关重要. \MBUnresolvedReference{sec:ch04-divergence-operator}{Section~3}和\MBUnresolvedReference{sec:ch04-curl-operator}{Section~4}讨论散度算子和旋度算子的特殊性质以及若干相关函数空间. \MBUnresolvedReference{sec:ch04-stokes-problem}{Section~5}分析带 Dirichlet 边界条件的 Stokes 问题, 其中包括 Stokes 算子的定义, 以及通过半群方法对非稳态 Stokes 问题的应用. 我们将 Stokes 方程椭圆正则性性质的完整证明推迟到\MBUnresolvedReference{sec:ch04-stokes-regularity}{Section~6}. \MBUnresolvedReference{sec:ch04-stress-boundary}{Section~7}, \MBUnresolvedReference{sec:ch04-interface-stokes}{Section~8}和\MBUnresolvedReference{sec:ch04-vorticity-boundary}{Section~9}讨论 Stokes 问题的一些非标准变体, 其中考虑不同类型的边界条件或界面条件.

\MBUnresolvedReference{chap:homogeneous-navier-stokes}{Chapter~V}研究稳态和非稳态 Navier--Stokes 方程. 对于非稳态问题, 我们证明全局弱解的存在性以及在二维情形下的唯一性, 还研究更光滑的解并讨论系统的抛物正则化性质. 随后, 对于稳态问题, 我们讨论齐次与非齐次边界数据情形下的存在性和唯一性问题, 并处理这些解的一些非常基本的稳定性问题.

在\MBUnresolvedReference{chap:nonhomogeneous-fluids}{Chapter~VI}中, 我们研究本书所考虑的最复杂模型, 即非均匀流体的(非稳态)不可压缩 Navier--Stokes 方程. 分析这一问题所需的第一个要素是研究输运方程的弱解. \MBUnresolvedReference{sec:ch06-transport-equation}{Section~1}以自成体系的方式介绍这类方程重整化解的 Di Perna--Lions 理论, 其中包括与相应迹问题有关的较新进展. 随后在\MBUnresolvedReference{sec:ch06-nonhomogeneous-navier-stokes}{Section~2}中利用这些结果, 证明完整 Navier--Stokes 系统所对应初边值问题弱解的全局存在性.

\MBUnresolvedReference{chap:boundary-conditions-modelling}{Chapter~VII}讨论不可压缩流数值模拟中出现的边界条件所涉及的两个不同问题. \MBUnresolvedReference{sec:ch07-outflow-boundary}{Section~1}处理出流边界条件问题. 当计算区域因实际原因小于真实物理区域时, 必须选择这类边界条件. 我们在此分析一种将出流边界上的法向应力与动量通量联系起来的非线性边界条件. \MBUnresolvedReference{sec:ch07-penalty-dirichlet}{Section~2}研究罚方法. 例如, 可以用这种方法在位于几何形状简单的计算区域内的障碍物边界上施加齐次 Dirichlet 边界条件. 通过描述一个适当的边界层, 我们证明当罚参数趋于 $0$ 时, 罚解收敛于精确解. 除了实际意义之外, 对这两个问题的分析还提供了引入若干数学方法的机会, 读者可能会发现这些方法在其他情形中也很有用.

\begin{flushright}
Marseille, 2012 年 8 月\\
Franck Boyer, Pierre Fabrie
\end{flushright}
